
\documentclass{article}
\usepackage{colortbl}
\usepackage{makecell}
\usepackage{multirow}
\usepackage{supertabular}

\begin{document}

\newcounter{utterance}

\centering \large Interaction Transcript for game `function\_detective', experiment `two\_numbers\_passive\_labeled\_pairs', episode 255 with Qwen3.6{-}35B{-}A3B{-}FP8{-}without{-}reasoning.
\vspace{24pt}

{ \footnotesize  \setcounter{utterance}{1}
\setlength{\tabcolsep}{0pt}
\begin{supertabular}{c@{$\;$}|p{.15\linewidth}@{}p{.15\linewidth}p{.15\linewidth}p{.15\linewidth}p{.15\linewidth}p{.15\linewidth}}
    \# & \multicolumn{2}{c}{Player} && \multicolumn{2}{c}{Game Master} \\
    \hline

    \theutterance \stepcounter{utterance}  
    & & & \multicolumn{4}{p{0.6\linewidth}}{
        \cellcolor[rgb]{0.9,0.9,0.9}{
            \makecell[{{p{\linewidth}}}]{
                \texttt{\tiny{[P1$\langle$GM]}}
                \texttt{I have a hidden Python function for you to figure out.} \\
\\ 
\texttt{Your job is to infer what it does and then submit a correct Python implementation.} \\
\\ 
\texttt{Function signature:} \\
\texttt{(a: int, b: int) {-}> int} \\
\\ 
\texttt{You have 15 turns total.} \\
\\ 
\texttt{I reveal candidate input{-}output pairs one at a time and tell you whether each one is correct.} \\
\\ 
\texttt{You must make your best SOLVE attempt no later than the second{-}to{-}last turn.} \\
\texttt{Do not wait until the very end.} \\
\\ 
\texttt{You can use these actions:} \\
\texttt{{-} JSON format described at the end of the prompt} \\
\texttt{{-} SOLVE: ```python ... ```} \\
\\ 
\texttt{Do not add import statements in your guess.} \\
\texttt{Anything you might need is already available, so just write the function body in hypothesis and return the answer.} \\
\\ 
\texttt{Preloaded examples:} \\
\texttt{1. ((({-}1000, 1), {-}2), True)} \\
\\ 
\texttt{Scoring:} \\
\texttt{{-} correctness is required for any non{-}zero score} \\
\texttt{{-} in interactive modes, solving with fewer information requests is better} \\
\\ 
\texttt{Input formatting reminders:} \\
\texttt{{-} strings must be quoted} \\
\texttt{{-} booleans must be True/False} \\
\texttt{{-} use commas between multiple values} \\
\texttt{{-} do not write any extra action tags in the same message} \\
\\ 
\texttt{Format example:} \\
\texttt{You write:} \\
\texttt{\{'current\_mode': <probing/confirming/don't know; choose whichever best describes your current state in the task>, 'input\_rationale': <brief description of why you chose this hypothesis>, 'hypothesis': <either the exact string "unknown" if you have no guess yet OR your current best guess as runnable Python matching the given signature, with a one{-}line docstring stating what it does. Never write placeholder code, if you cannot commit to a guess, use "unknown". Match the given signature.} \\
\texttt{e.g.} \\
\texttt{def f(x):} \\
\texttt{	"""Returns the absolute value of x."""} \\
\texttt{	return x if x >= 0 else {-}x} \\
\\ 
\texttt{Use "unknown" if you have no guess.>, 'confidence\_score': <score between 1{-}100. 100 means highest confidence in your hypothesis>\}} \\
\\ 
\texttt{Output only one valid json per turn.} \\
            }
        }
    }
     \\ \\

    \theutterance \stepcounter{utterance}  
    & \multicolumn{4}{p{0.6\linewidth}}{
        \cellcolor[rgb]{0.561,0.737,0.561}{
            \makecell[{{p{\linewidth}}}]{
                \texttt{\tiny{[P1$\rangle$GM]}}
                \texttt{\{'current\_mode': 'probing', 'input\_rationale': 'I have one preloaded example: f({-}1000, 1) = True. This is strange because the signature says it returns an int, but the example shows True (which is a bool, but in Python True == 1). Let me probe with more inputs to understand the behavior. I'll try simple cases like (0, 0), (1, 0), (0, 1), (1, 1) to see if it's a comparison, arithmetic, or bitwise operation.', 'hypothesis': 'unknown', 'confidence\_score': 10\}} \\
            }
        }
    }
    & & \\ \\

    \theutterance \stepcounter{utterance}  
    & & & \multicolumn{4}{p{0.6\linewidth}}{
        \cellcolor[rgb]{0.9,0.9,0.9}{
            \makecell[{{p{\linewidth}}}]{
                \texttt{\tiny{[P1$\langle$GM]}}
                \texttt{Game Violation: ParseError: Response is not valid JSON.} \\
\texttt{Output ONLY:} \\
\texttt{One JSON with confidence\_score, current\_mode, input\_rationale, and hypothesis keys ...} \\
\texttt{OR} \\
\texttt{SOLVE: ```python ... ```} \\
            }
        }
    }
     \\ \\

    \theutterance \stepcounter{utterance}  
    & \multicolumn{4}{p{0.6\linewidth}}{
        \cellcolor[rgb]{0.561,0.737,0.561}{
            \makecell[{{p{\linewidth}}}]{
                \texttt{\tiny{[P1$\rangle$GM]}}
                \texttt{\{} \\
\texttt{  "current\_mode": "probing",} \\
\texttt{  "input\_rationale": "I have one preloaded example: f({-}1000, 1) = True. Since the return type is int, True likely means 1. I need to probe more to understand the function. Let me try simple cases like (0, 0), (1, 0), (0, 1), (1, 1) to see if it's a comparison, arithmetic, or bitwise operation.",} \\
\texttt{  "hypothesis": "unknown",} \\
\texttt{  "confidence\_score": 10} \\
\texttt{\}} \\
            }
        }
    }
    & & \\ \\

    \theutterance \stepcounter{utterance}  
    & & & \multicolumn{4}{p{0.6\linewidth}}{
        \cellcolor[rgb]{0.9,0.9,0.9}{
            \makecell[{{p{\linewidth}}}]{
                \texttt{\tiny{[P1$\langle$GM]}}
                \texttt{OUTPUT: (((0, {-}2), 1), False)} \\
            }
        }
    }
     \\ \\

    \theutterance \stepcounter{utterance}  
    & \multicolumn{4}{p{0.6\linewidth}}{
        \cellcolor[rgb]{0.561,0.737,0.561}{
            \makecell[{{p{\linewidth}}}]{
                \texttt{\tiny{[P1$\rangle$GM]}}
                \texttt{\{} \\
\texttt{  "current\_mode": "probing",} \\
\texttt{  "input\_rationale": "I have two data points: f({-}1000, 1) = True (1) and f(0, {-}2) = False (0). The first argument is large negative, second is positive {-}> True. Second: first is 0, second is negative {-}> False. This suggests the function might be checking if the first argument is greater than the second, or some other comparison. Let me test f(1, 0) and f(0, 1) to distinguish between > and >= or other operations.",} \\
\texttt{  "hypothesis": "unknown",} \\
\texttt{  "confidence\_score": 20} \\
\texttt{\}} \\
            }
        }
    }
    & & \\ \\

    \theutterance \stepcounter{utterance}  
    & & & \multicolumn{4}{p{0.6\linewidth}}{
        \cellcolor[rgb]{0.9,0.9,0.9}{
            \makecell[{{p{\linewidth}}}]{
                \texttt{\tiny{[P1$\langle$GM]}}
                \texttt{OUTPUT: (((2, 2), 4), True)} \\
            }
        }
    }
     \\ \\

    \theutterance \stepcounter{utterance}  
    & \multicolumn{4}{p{0.6\linewidth}}{
        \cellcolor[rgb]{0.561,0.737,0.561}{
            \makecell[{{p{\linewidth}}}]{
                \texttt{\tiny{[P1$\rangle$GM]}}
                \texttt{\{} \\
\texttt{  "current\_mode": "probing",} \\
\texttt{  "input\_rationale": "I have three data points: f({-}1000, 1) = True, f(0, {-}2) = False, f(2, 2) = True. The function returns True (1) or False (0). Let's analyze: {-}1000 > 1 is False, but result is True. 0 > {-}2 is True, but result is False. 2 > 2 is False, but result is True. So it's not simple >. Let's check if it's checking if the first argument is less than the second: {-}1000 < 1 is True (matches), 0 < {-}2 is False (matches), 2 < 2 is False (but result is True, so doesn't match). Let's check if it's checking if the first argument is not equal to the second: {-}1000 != 1 is True (matches), 0 != {-}2 is True (but result is False, so doesn't match). Let me try more inputs to narrow it down. Perhaps it's related to parity or some other property. Let me test f(1, 0) and f(0, 1).",} \\
\texttt{  "hypothesis": "unknown",} \\
\texttt{  "confidence\_score": 25} \\
\texttt{\}} \\
            }
        }
    }
    & & \\ \\

    \theutterance \stepcounter{utterance}  
    & & & \multicolumn{4}{p{0.6\linewidth}}{
        \cellcolor[rgb]{0.9,0.9,0.9}{
            \makecell[{{p{\linewidth}}}]{
                \texttt{\tiny{[P1$\langle$GM]}}
                \texttt{OUTPUT: ((({-}1, 1), {-}1), False)} \\
            }
        }
    }
     \\ \\

    \theutterance \stepcounter{utterance}  
    & \multicolumn{4}{p{0.6\linewidth}}{
        \cellcolor[rgb]{0.561,0.737,0.561}{
            \makecell[{{p{\linewidth}}}]{
                \texttt{\tiny{[P1$\rangle$GM]}}
                \texttt{\{} \\
\texttt{  "current\_mode": "probing",} \\
\texttt{  "input\_rationale": "I have four data points: f({-}1000, 1) = True, f(0, {-}2) = False, f(2, 2) = True, f({-}1, 1) = False. Let's look at the pattern more carefully. The return value is 1 (True) or 0 (False). Let's check if it's related to the sign of the first argument or some bitwise operation. {-}1000 is negative, result True. 0 is non{-}negative, result False. 2 is positive, result True. {-}1 is negative, result False. This doesn't seem to be about the sign of the first argument alone. Let me check if it's about the second argument. 1 is positive, True. {-}2 is negative, False. 2 is positive, True. 1 is positive, False. So it's not just the sign of the second argument either. Let me try f(1, 1) and f({-}1, {-}1) to see if there's a pattern with equal arguments or same signs.",} \\
\texttt{  "hypothesis": "unknown",} \\
\texttt{  "confidence\_score": 30} \\
\texttt{\}} \\
            }
        }
    }
    & & \\ \\

    \theutterance \stepcounter{utterance}  
    & & & \multicolumn{4}{p{0.6\linewidth}}{
        \cellcolor[rgb]{0.9,0.9,0.9}{
            \makecell[{{p{\linewidth}}}]{
                \texttt{\tiny{[P1$\langle$GM]}}
                \texttt{OUTPUT: ((({-}1000, {-}99), 198), True)} \\
            }
        }
    }
     \\ \\

    \theutterance \stepcounter{utterance}  
    & \multicolumn{4}{p{0.6\linewidth}}{
        \cellcolor[rgb]{0.561,0.737,0.561}{
            \makecell[{{p{\linewidth}}}]{
                \texttt{\tiny{[P1$\rangle$GM]}}
                \texttt{\{} \\
\texttt{  "current\_mode": "probing",} \\
\texttt{  "input\_rationale": "I have five data points: f({-}1000, 1) = True, f(0, {-}2) = False, f(2, 2) = True, f({-}1, 1) = False, f({-}1000, {-}99) = True. Let's look at the relationship between a and b. Maybe it's checking if a is divisible by b? {-}1000 \% 1 == 0 {-}> True. 0 \% {-}2 == 0 {-}> True, but result is False. So not divisibility. Maybe it's checking if a + b is even? {-}1000 + 1 = {-}999 (odd) {-}> True. 0 + ({-}2) = {-}2 (even) {-}> False. 2 + 2 = 4 (even) {-}> True. Doesn't match. Maybe it's checking if a * b is positive? {-}1000 * 1 = {-}1000 (negative) {-}> True. 0 * {-}2 = 0 {-}> False. 2 * 2 = 4 (positive) {-}> True. {-}1 * 1 = {-}1 (negative) {-}> False. {-}1000 * {-}99 = 99000 (positive) {-}> True. This doesn't seem consistent either} \\
            }
        }
    }
    & & \\ \\

    \theutterance \stepcounter{utterance}  
    & & & \multicolumn{4}{p{0.6\linewidth}}{
        \cellcolor[rgb]{0.9,0.9,0.9}{
            \makecell[{{p{\linewidth}}}]{
                \texttt{\tiny{[P1$\langle$GM]}}
                \texttt{Game Violation: ParseError: Response is not valid JSON.} \\
\texttt{Output ONLY:} \\
\texttt{One JSON with confidence\_score, current\_mode, input\_rationale, and hypothesis keys ...} \\
\texttt{OR} \\
\texttt{SOLVE: ```python ... ```} \\
            }
        }
    }
     \\ \\

    \theutterance \stepcounter{utterance}  
    & \multicolumn{4}{p{0.6\linewidth}}{
        \cellcolor[rgb]{0.561,0.737,0.561}{
            \makecell[{{p{\linewidth}}}]{
                \texttt{\tiny{[P1$\rangle$GM]}}
                \texttt{\{} \\
\texttt{  "current\_mode": "probing",} \\
\texttt{  "input\_rationale": "I have five data points: f({-}1000, 1) = True, f(0, {-}2) = False, f(2, 2) = True, f({-}1, 1) = False, f({-}1000, {-}99) = True. Let's look at the relationship between a and b. Maybe it's checking if a is divisible by b? {-}1000 \% 1 == 0 {-}> True. 0 \% {-}2 == 0 {-}> True, but result is False. So not divisibility. Maybe it's checking if a + b is even? {-}1000 + 1 = {-}999 (odd) {-}> True. 0 + ({-}2) = {-}2 (even) {-}> False. 2 + 2 = 4 (even) {-}> True. Doesn't match. Maybe it's checking if a * b is positive? {-}1000 * 1 = {-}1000 (negative) {-}> True. 0 * {-}2 = 0 {-}> False. 2 * 2 = 4 (positive) {-}> True. {-}1 * 1 = {-}1 (negative) {-}> False. {-}1000 * {-}99 = 99000 (positive) {-}> True. This doesn't seem consistent either} \\
            }
        }
    }
    & & \\ \\

    \theutterance \stepcounter{utterance}  
    & & & \multicolumn{4}{p{0.6\linewidth}}{
        \cellcolor[rgb]{0.9,0.9,0.9}{
            \makecell[{{p{\linewidth}}}]{
                \texttt{\tiny{[P1$\langle$GM]}}
                \texttt{Game Violation: ParseError: Response is not valid JSON.} \\
\texttt{Output ONLY:} \\
\texttt{One JSON with confidence\_score, current\_mode, input\_rationale, and hypothesis keys ...} \\
\texttt{OR} \\
\texttt{SOLVE: ```python ... ```} \\
            }
        }
    }
     \\ \\

    \theutterance \stepcounter{utterance}  
    & \multicolumn{4}{p{0.6\linewidth}}{
        \cellcolor[rgb]{0.561,0.737,0.561}{
            \makecell[{{p{\linewidth}}}]{
                \texttt{\tiny{[P1$\rangle$GM]}}
                \texttt{\{} \\
\texttt{  "current\_mode": "probing",} \\
\texttt{  "input\_rationale": "I have five data points: f({-}1000, 1) = True, f(0, {-}2) = False, f(2, 2) = True, f({-}1, 1) = False, f({-}1000, {-}99) = True. Let's look at the relationship between a and b. Maybe it's checking if a is divisible by b? {-}1000 \% 1 == 0 {-}> True. 0 \% {-}2 == 0 {-}> True, but result is False. So not divisibility. Maybe it's checking if a + b is even? {-}1000 + 1 = {-}999 (odd) {-}> True. 0 + ({-}2) = {-}2 (even) {-}> False. 2 + 2 = 4 (even) {-}> True. Doesn't match. Maybe it's checking if a * b is positive? {-}1000 * 1 = {-}1000 (negative) {-}> True. 0 * {-}2 = 0 {-}> False. 2 * 2 = 4 (positive) {-}> True. {-}1 * 1 = {-}1 (negative) {-}> False. {-}1000 * {-}99 = 99000 (positive) {-}> True. This doesn't seem consistent either} \\
            }
        }
    }
    & & \\ \\

    \theutterance \stepcounter{utterance}  
    & & & \multicolumn{4}{p{0.6\linewidth}}{
        \cellcolor[rgb]{0.9,0.9,0.9}{
            \makecell[{{p{\linewidth}}}]{
                \texttt{\tiny{[P1$\langle$GM]}}
                \texttt{Game Violation: ParseError: Response is not valid JSON.} \\
\texttt{Output ONLY:} \\
\texttt{One JSON with confidence\_score, current\_mode, input\_rationale, and hypothesis keys ...} \\
\texttt{OR} \\
\texttt{SOLVE: ```python ... ```} \\
            }
        }
    }
     \\ \\

    \theutterance \stepcounter{utterance}  
    & \multicolumn{4}{p{0.6\linewidth}}{
        \cellcolor[rgb]{0.561,0.737,0.561}{
            \makecell[{{p{\linewidth}}}]{
                \texttt{\tiny{[P1$\rangle$GM]}}
                \texttt{\{} \\
\texttt{  "current\_mode": "probing",} \\
\texttt{  "input\_rationale": "I have five data points: f({-}1000, 1) = True, f(0, {-}2) = False, f(2, 2) = True, f({-}1, 1) = False, f({-}1000, {-}99) = True. Let's look at the relationship between a and b. Maybe it's checking if a is divisible by b? {-}1000 \% 1 == 0 {-}> True. 0 \% {-}2 == 0 {-}> True, but result is False. So not divisibility. Maybe it's checking if a + b is even? {-}1000 + 1 = {-}999 (odd) {-}> True. 0 + ({-}2) = {-}2 (even) {-}> False. 2 + 2 = 4 (even) {-}> True. Doesn't match. Maybe it's checking if a * b is positive? {-}1000 * 1 = {-}1000 (negative) {-}> True. 0 * {-}2 = 0 {-}> False. 2 * 2 = 4 (positive) {-}> True. {-}1 * 1 = {-}1 (negative) {-}> False. {-}1000 * {-}99 = 99000 (positive) {-}> True. This doesn't seem consistent either} \\
            }
        }
    }
    & & \\ \\

    \theutterance \stepcounter{utterance}  
    & & & \multicolumn{4}{p{0.6\linewidth}}{
        \cellcolor[rgb]{0.9,0.9,0.9}{
            \makecell[{{p{\linewidth}}}]{
                \texttt{\tiny{[P1$\langle$GM]}}
                \texttt{Game Violation: ParseError: Response is not valid JSON.} \\
\texttt{Output ONLY:} \\
\texttt{One JSON with confidence\_score, current\_mode, input\_rationale, and hypothesis keys ...} \\
\texttt{OR} \\
\texttt{SOLVE: ```python ... ```} \\
            }
        }
    }
     \\ \\

    \theutterance \stepcounter{utterance}  
    & \multicolumn{4}{p{0.6\linewidth}}{
        \cellcolor[rgb]{0.561,0.737,0.561}{
            \makecell[{{p{\linewidth}}}]{
                \texttt{\tiny{[P1$\rangle$GM]}}
                \texttt{\{} \\
\texttt{  "current\_mode": "probing",} \\
\texttt{  "input\_rationale": "I have five data points: f({-}1000, 1) = True, f(0, {-}2) = False, f(2, 2) = True, f({-}1, 1) = False, f({-}1000, {-}99) = True. Let's look at the relationship between a and b. Maybe it's checking if a is divisible by b? {-}1000 \% 1 == 0 {-}> True. 0 \% {-}2 == 0 {-}> True, but result is False. So not divisibility. Maybe it's checking if a + b is even? {-}1000 + 1 = {-}999 (odd) {-}> True. 0 + ({-}2) = {-}2 (even) {-}> False. 2 + 2 = 4 (even) {-}> True. Doesn't match. Maybe it's checking if a * b is positive? {-}1000 * 1 = {-}1000 (negative) {-}> True. 0 * {-}2 = 0 {-}> False. 2 * 2 = 4 (positive) {-}> True. {-}1 * 1 = {-}1 (negative) {-}> False. {-}1000 * {-}99 = 99000 (positive) {-}> True. This doesn't seem consistent either} \\
            }
        }
    }
    & & \\ \\

    \theutterance \stepcounter{utterance}  
    & & & \multicolumn{4}{p{0.6\linewidth}}{
        \cellcolor[rgb]{0.9,0.9,0.9}{
            \makecell[{{p{\linewidth}}}]{
                \texttt{\tiny{[P1$\langle$GM]}}
                \texttt{Game Violation: ParseError: Response is not valid JSON.} \\
\texttt{Output ONLY:} \\
\texttt{One JSON with confidence\_score, current\_mode, input\_rationale, and hypothesis keys ...} \\
\texttt{OR} \\
\texttt{SOLVE: ```python ... ```} \\
            }
        }
    }
     \\ \\

    \theutterance \stepcounter{utterance}  
    & \multicolumn{4}{p{0.6\linewidth}}{
        \cellcolor[rgb]{0.561,0.737,0.561}{
            \makecell[{{p{\linewidth}}}]{
                \texttt{\tiny{[P1$\rangle$GM]}}
                \texttt{\{} \\
\texttt{  "current\_mode": "probing",} \\
\texttt{  "input\_rationale": "I have five data points: f({-}1000, 1) = True, f(0, {-}2) = False, f(2, 2) = True, f({-}1, 1) = False, f({-}1000, {-}99) = True. Let's look at the relationship between a and b. Maybe it's checking if a is divisible by b? {-}1000 \% 1 == 0 {-}> True. 0 \% {-}2 == 0 {-}> True, but result is False. So not divisibility. Maybe it's checking if a + b is even? {-}1000 + 1 = {-}999 (odd) {-}> True. 0 + ({-}2) = {-}2 (even) {-}> False. 2 + 2 = 4 (even) {-}> True. Doesn't match. Maybe it's checking if a * b is positive? {-}1000 * 1 = {-}1000 (negative) {-}> True. 0 * {-}2 = 0 {-}> False. 2 * 2 = 4 (positive) {-}> True. {-}1 * 1 = {-}1 (negative) {-}> False. {-}1000 * {-}99 = 99000 (positive) {-}> True. This doesn't seem consistent either} \\
            }
        }
    }
    & & \\ \\

    \theutterance \stepcounter{utterance}  
    & & & \multicolumn{4}{p{0.6\linewidth}}{
        \cellcolor[rgb]{0.9,0.9,0.9}{
            \makecell[{{p{\linewidth}}}]{
                \texttt{\tiny{[P1$\langle$GM]}}
                \texttt{Game Violation: ParseError: Response is not valid JSON.} \\
\texttt{Output ONLY:} \\
\texttt{One JSON with confidence\_score, current\_mode, input\_rationale, and hypothesis keys ...} \\
\texttt{OR} \\
\texttt{SOLVE: ```python ... ```} \\
            }
        }
    }
     \\ \\

    \theutterance \stepcounter{utterance}  
    & \multicolumn{4}{p{0.6\linewidth}}{
        \cellcolor[rgb]{0.561,0.737,0.561}{
            \makecell[{{p{\linewidth}}}]{
                \texttt{\tiny{[P1$\rangle$GM]}}
                \texttt{\{} \\
\texttt{  "current\_mode": "probing",} \\
\texttt{  "input\_rationale": "I have five data points: f({-}1000, 1) = True, f(0, {-}2) = False, f(2, 2) = True, f({-}1, 1) = False, f({-}1000, {-}99) = True. Let's look at the relationship between a and b. Maybe it's checking if a is divisible by b? {-}1000 \% 1 == 0 {-}> True. 0 \% {-}2 == 0 {-}> True, but result is False. So not divisibility. Maybe it's checking if a + b is even? {-}1000 + 1 = {-}999 (odd) {-}> True. 0 + ({-}2) = {-}2 (even) {-}> False. 2 + 2 = 4 (even) {-}> True. Doesn't match. Maybe it's checking if a * b is positive? {-}1000 * 1 = {-}1000 (negative) {-}> True. 0 * {-}2 = 0 {-}> False. 2 * 2 = 4 (positive) {-}> True. {-}1 * 1 = {-}1 (negative) {-}> False. {-}1000 * {-}99 = 99000 (positive) {-}> True. This doesn't seem consistent either} \\
            }
        }
    }
    & & \\ \\

    \theutterance \stepcounter{utterance}  
    & & & \multicolumn{4}{p{0.6\linewidth}}{
        \cellcolor[rgb]{0.9,0.9,0.9}{
            \makecell[{{p{\linewidth}}}]{
                \texttt{\tiny{[P1$\langle$GM]}}
                \texttt{Game Violation: ParseError: Response is not valid JSON.} \\
\texttt{Output ONLY:} \\
\texttt{One JSON with confidence\_score, current\_mode, input\_rationale, and hypothesis keys ...} \\
\texttt{OR} \\
\texttt{SOLVE: ```python ... ```} \\
            }
        }
    }
     \\ \\

    \theutterance \stepcounter{utterance}  
    & \multicolumn{4}{p{0.6\linewidth}}{
        \cellcolor[rgb]{0.561,0.737,0.561}{
            \makecell[{{p{\linewidth}}}]{
                \texttt{\tiny{[P1$\rangle$GM]}}
                \texttt{\{} \\
\texttt{  "current\_mode": "probing",} \\
\texttt{  "input\_rationale": "I have five data points: f({-}1000, 1) = True, f(0, {-}2) = False, f(2, 2) = True, f({-}1, 1) = False, f({-}1000, {-}99) = True. Let's look at the relationship between a and b. Maybe it's checking if a is divisible by b? {-}1000 \% 1 == 0 {-}> True. 0 \% {-}2 == 0 {-}> True, but result is False. So not divisibility. Maybe it's checking if a + b is even? {-}1000 + 1 = {-}999 (odd) {-}> True. 0 + ({-}2) = {-}2 (even) {-}> False. 2 + 2 = 4 (even) {-}> True. Doesn't match. Maybe it's checking if a * b is positive? {-}1000 * 1 = {-}1000 (negative) {-}> True. 0 * {-}2 = 0 {-}> False. 2 * 2 = 4 (positive) {-}> True. {-}1 * 1 = {-}1 (negative) {-}> False. {-}1000 * {-}99 = 99000 (positive) {-}> True. This doesn't seem consistent either} \\
            }
        }
    }
    & & \\ \\

    \theutterance \stepcounter{utterance}  
    & & & \multicolumn{4}{p{0.6\linewidth}}{
        \cellcolor[rgb]{0.9,0.9,0.9}{
            \makecell[{{p{\linewidth}}}]{
                \texttt{\tiny{[P1$\langle$GM]}}
                \texttt{Game Violation: ParseError: Response is not valid JSON.} \\
\texttt{Output ONLY:} \\
\texttt{One JSON with confidence\_score, current\_mode, input\_rationale, and hypothesis keys ...} \\
\texttt{OR} \\
\texttt{SOLVE: ```python ... ```} \\
            }
        }
    }
     \\ \\

    \theutterance \stepcounter{utterance}  
    & \multicolumn{4}{p{0.6\linewidth}}{
        \cellcolor[rgb]{0.561,0.737,0.561}{
            \makecell[{{p{\linewidth}}}]{
                \texttt{\tiny{[P1$\rangle$GM]}}
                \texttt{\{} \\
\texttt{  "current\_mode": "probing",} \\
\texttt{  "input\_rationale": "I have five data points: f({-}1000, 1) = True, f(0, {-}2) = False, f(2, 2) = True, f({-}1, 1) = False, f({-}1000, {-}99) = True. Let's look at the relationship between a and b. Maybe it's checking if a is divisible by b? {-}1000 \% 1 == 0 {-}> True. 0 \% {-}2 == 0 {-}> True, but result is False. So not divisibility. Maybe it's checking if a + b is even? {-}1000 + 1 = {-}999 (odd) {-}> True. 0 + ({-}2) = {-}2 (even) {-}> False. 2 + 2 = 4 (even) {-}> True. Doesn't match. Maybe it's checking if a * b is positive? {-}1000 * 1 = {-}1000 (negative) {-}> True. 0 * {-}2 = 0 {-}> False. 2 * 2 = 4 (positive) {-}> True. {-}1 * 1 = {-}1 (negative) {-}> False. {-}1000 * {-}99 = 99000 (positive) {-}> True. This doesn't seem consistent either} \\
            }
        }
    }
    & & \\ \\

    \theutterance \stepcounter{utterance}  
    & & & \multicolumn{4}{p{0.6\linewidth}}{
        \cellcolor[rgb]{0.9,0.9,0.9}{
            \makecell[{{p{\linewidth}}}]{
                \texttt{\tiny{[P1$\langle$GM]}}
                \texttt{Game Violation: ParseError: Response is not valid JSON.} \\
\texttt{Output ONLY:} \\
\texttt{One JSON with confidence\_score, current\_mode, input\_rationale, and hypothesis keys ...} \\
\texttt{OR} \\
\texttt{SOLVE: ```python ... ```} \\
            }
        }
    }
     \\ \\

    \theutterance \stepcounter{utterance}  
    & \multicolumn{4}{p{0.6\linewidth}}{
        \cellcolor[rgb]{0.561,0.737,0.561}{
            \makecell[{{p{\linewidth}}}]{
                \texttt{\tiny{[P1$\rangle$GM]}}
                \texttt{\{} \\
\texttt{  "current\_mode": "probing",} \\
\texttt{  "input\_rationale": "I have five data points: f({-}1000, 1) = True, f(0, {-}2) = False, f(2, 2) = True, f({-}1, 1) = False, f({-}1000, {-}99) = True. Let's look at the relationship between a and b. Maybe it's checking if a is divisible by b? {-}1000 \% 1 == 0 {-}> True. 0 \% {-}2 == 0 {-}> True, but result is False. So not divisibility. Maybe it's checking if a + b is even? {-}1000 + 1 = {-}999 (odd) {-}> True. 0 + ({-}2) = {-}2 (even) {-}> False. 2 + 2 = 4 (even) {-}> True. Doesn't match. Maybe it's checking if a * b is positive? {-}1000 * 1 = {-}1000 (negative) {-}> True. 0 * {-}2 = 0 {-}> False. 2 * 2 = 4 (positive) {-}> True. {-}1 * 1 = {-}1 (negative) {-}> False. {-}1000 * {-}99 = 99000 (positive) {-}> True. This doesn't seem consistent either} \\
            }
        }
    }
    & & \\ \\

    \theutterance \stepcounter{utterance}  
    & & & \multicolumn{2}{p{0.3\linewidth}}{
        \cellcolor[rgb]{0.9,0.9,0.9}{
            \makecell[{{p{\linewidth}}}]{
                \texttt{\tiny{[GM$|$GM]}}
                \texttt{Game Verdict: LOSS (turn limit reached)} \\
            }
        }
    }
    & & \\ \\

\end{supertabular}
}

\end{document}
